\scanpage{45}
\mspage{2}\margin{44}
\deleted{donc} \inserted{ou encore,} qu'on a un élément
\deleted{privilégié \unsure{d'ordre}} \inserted{d'ordre 2}
\[
(7)\qquad a_0\deleted{\unsure{a}}\in D^+
\qquad a_0:x\longmapsto-x
\]
qui est en fait le seul élément \inserted{d'ordre 2} $D^+$.
Si à tout $j\in\widetilde{\unsure{I|J}}$ on associe l'unique
\deleted{\unsure{réflexe}} élément $r_j\in D$ (une réflexion)
\deleted{donc} qui soit $\ne1$ et telle que
\[
(8)\qquad r_j(\widetilde R_j)=\widetilde R_j
\]
on aura
\[
(9)\qquad r_j=r_{a_0(j)},
\]
\deleted{et} donc il n'y a que $n=m/2$ réflexions $r_j$ distinctes,
qui correspondent à la moitié des éléments de la classe \inserted{non neutre}
$D^-$ de $D$ modulo $D^+$. Les autres réflexions sont celles associées aux
\uline{côtés} du polygone envisagé.

Il y a sur $V$ une unique forme quadratique invariante par $D$
(ou seulement par $D^+$), et égale à 1 sur les $\widetilde R_j$.
On considère $V$ comme espace \inserted{vect.} euclidien, grâce à cette forme.

\margin{\uline{NB} \deleted{$\unsure{a_0}$} \inserted{\unsure{Le centre} de}
$D$ est réduit à $\{1,a_0\}$, $m$ pair signifie \unsure{donc que}
le centre de $D$ est $\ne\{1\}$.}
\margin{\deleted{$\unsure{a_0}$ \illegible{2} \unsure{qui fixe}
$\unsure{aR}$ \unsure{des vecteurs positifs}} \deleted{\unsure{point}}
d'ordre 2}
